% !TeX root = ../queens.tex
\section{Fast generation with little memory}\label{sec:generation}

The local description also gives a practical way to compute the
queens. The resulting program generates $q_0,\dots,q_N$ with linear
total work and logarithmic working memory: it keeps neither the
board nor the growing queen word. It rests on two ideas. First, the
calculation for column $n$ reads the queen word only far behind
column $n$, so the symbols it needs can be regenerated by a second
copy of the same calculation running behind the first
(Sections~\ref{sec:generator-idea} and~\ref{sec:generator-records}).
Second, along the actual process the calculation passes through only
finitely many records, so it can be compiled into a fixed table that
places several queens per lookup, with no attack tests at run time
(Section~\ref{sec:block-transducer}).
Section~\ref{sec:generator-cost} adds buffering between the copies and
proves the bounds, and Section~\ref{sec:generation-benchmarks}
compares the C implementation with bit-packed Knuth.

\subsection{Regenerating the earlier word}\label{sec:generator-idea}

By Lemma~\ref{lem:local-exactness}, the actual branch of the
calculation for column $n$ needs the queen word only at indices up
to $m+6$, where $m$ is the least unused row. On the board,
$n=m+\kappa+z$ with $\kappa=U(m-1)$ and $z$ bounded, and
Proposition~\ref{prop:stability}, with $C=4$ from
Lemma~\ref{lem:estimates}, gives $U(m-1)=m/\phi+O(1)$. Since
$1+1/\phi=\phi$, this means $n=\phi m+O(1)$: the calculation writes
$\sigma_n$ but reads the word only up to about index $n/\phi\approx
0.618\,n$.

The earlier symbols can therefore be regenerated instead of stored.
Start a copy of the calculation at the board before column $30$
(Section~\ref{app:seed}). Its records suffice to produce
$\sigma_{30},\dots,\sigma_{47}$; after that, it requests
$\sigma_{30},\sigma_{31},\dots$ in order. The copy has produced these
symbols itself, but it has not kept them. A second copy, started at
the same board, produces the same word and can answer these requests.
When the second copy needs input in turn, a third copy supplies it,
and so on, forming a chain of copies.
Every copy produces the same word from $\sigma_{30}$ onward, at the
pace required by the copy to its left. Because its inputs are the
actual symbols, each copy follows the actual branch of
Lemma~\ref{lem:local-exactness}: there is nothing to choose, and the
history graph is not consulted.

When the outermost copy reaches column $N$, the second copy works
near column $N/\phi$, the third near $N/\phi^2$, and so on, until a
copy is still within its first eighteen columns. For $N=10^6$, with
symbols passed one at a time, the second, third, and fourth copies
stand before columns $618035$, $381967$, and $236072$, and the chain
has $22$ copies, the last before column $48$
(Figure~\ref{fig:recursive-generation}). The number of copies thus grows logarithmically
with $N$, while their combined work is about
\[
 N+\frac{N}{\phi}+\frac{N}{\phi^2}+\cdots=\phi^2N\approx2.618\,N
\]
queen placements; for $N=10^6$ the copies carried out $2617400$
placements in total. Section~\ref{sec:generator-cost} makes this
estimate precise.

\begin{figure}[!htbp]
\centering
% Positions of the copies, in units of 10^5 columns, from the simulation
% described in Section 7.1 (symbols passed one at a time, N = 10^6).
\begin{tikzpicture}[x=1.35cm,y=1cm,>=stealth]
 \draw[->,black!70] (0,0) -- (10.45,0) node[right,font=\footnotesize] {column};
 \foreach \c in {10,6.18035,3.81967,2.36072,1.45904,0.90176,0.55734,
                 0.34448,0.21290,0.13159,0.08137,0.05030,0.03109,0.01923,
                 0.01190,0.00739,0.00457,0.00284,0.00177,0.00111,0.00070,0.00048}
  \draw[black!80,very thin] (\c,-.09) -- (\c,.09);
 \foreach \a/\b in {10/6.18035,6.18035/3.81967,3.81967/2.36072,
                    2.36072/1.45904,1.45904/0.90176,0.90176/0.55734,
                    0.55734/0.34448,0.34448/0.21290}
  \draw[->,queensUpper,semithick] (\a,.12) to[bend right=55] (\b,.12);
 \node[font=\footnotesize,text=queensUpper,above] at (8.1,.95) {reads};
 \foreach \c/\k/\v in {10/1/{1\,000\,000},6.18035/2/{618\,035},
                       3.81967/3/{381\,967},2.36072/4/{236\,072}}
  \node[below=3pt,align=center,font=\footnotesize] at (\c,0)
   {copy~\k\\$\v$};
 \node[below=3pt,font=\footnotesize] at (0,0) {$0$};
\end{tikzpicture}
\caption{Where the copies work when the outermost copy reaches column
$N=10^6$, with symbols passed one at a time. Each tick marks the next
column of one copy. Each arc leads from where a copy writes to where
it reads, which is where the next copy writes. The ticks crowd towards
the left and end with the $22$nd copy, before column $48$.}
\label{fig:recursive-generation}
\end{figure}


\subsection{The generator, one symbol at a time}\label{sec:generator-records}

Because a copy receives the actual symbols, it needs less than the
state of Section~\ref{sec:local}. The history graph and the output
history $H_{\mathrm{out}}$ serve only to constrain the possible inputs
and to check the outputs during verification. With the actual inputs there are
no choices to enumerate, so a copy keeps neither.

We can also compute the next row bit directly. By
\eqref{eq:source-offsets} and $n=m+\kappa+z$, an upper queen in column
$m+h$ occupies row $n$ exactly when
\begin{equation}\label{eq:generator-row-bit}
 u_{m+h}=1\quad\hbox{and}\quad h+\Gamma(h)=z.
\end{equation}
Thus $b_n$ is $1$ if this equality holds for a stored upper column,
and $0$ otherwise. This replaces the two counts $\J(z)$ and $\J(z-1)$
in \eqref{eq:output}.

To gather the needed symbols in one pass, first extend $Q$ to length
at least $\max(1,z,w+1)$, keeping it if it is already longer. This
covers every candidate row and every upper column needed for either
test: for $r\le w$, the request length
$\max\{r+1,\lfloor(z+r)/2\rfloor+1\}$ in Section~\ref{app:output} is at
most $\max(1,z,w+1)$. The latter is at most five, so these
preliminary requests end by $m+4$. The subsequent row search reads
further symbols as before, at most through $m+6$. Reading a few
symbols earlier does not change the queen choice or the row and
diagonal records.

Only four column bits before $m$ are needed. Indeed, for
$h\le-5$, we have $\Gamma(h)\le0$ and hence
\[
 h+\Gamma(h)\le-5<z,\qquad
 2h+\Gamma(h)\le-10<z+r
\]
for every lower candidate $r\ge0$, using $z\ge-4$. Such a column
cannot satisfy either \eqref{eq:generator-row-bit} or the
upper-antidiagonal test \eqref{eq:upper-sum}. The row tests on the
lower candidates read their bits directly from $Q$, so past row bits
are unnecessary.

The records of a copy can consequently be just
\begin{equation}\label{eq:generator-records}
 \bigl(w,z,R,D,A,(u_{m-4},u_{m-3},u_{m-2},u_{m-1}),Q\bigr).
\end{equation}
The field $z$ is updated by \eqref{eq:update}, whose sum uses the
symbols removed from $Q$; append their column bits to the four
stored bits, retaining only the last four. All other updates are
unchanged. In particular, $Q$ still stores both bits of each symbol.
Its length stays at most eleven: it starts at eleven, is extended
only to a required length of at most seven, and then loses its first
$\mu$ symbols. The starting values, before column $30$, are
\[
 \begin{gathered}
 w=0,\quad z=-1,\quad R=A=\varnothing,\quad D=\{1\},\\
 (u_{15},u_{16},u_{17},u_{18})=(1,1,1,1),\qquad
 Q=\mathtt{00322301212}.
 \end{gathered}
\]

A copy consists of these records and a list of the symbols it has
computed but not yet handed on. Algorithm~\ref{alg:next-symbol}
hands on its next symbol, asking its own input copy for symbols as
needed. The outermost copy runs the same loop but recovers the rows
instead of handing on symbols. It keeps the column $n$, the row
reference $m$, and the upper count $U(n-1)$ as ordinary integers: a
lower choice with offset $r$ gives $q_n=m+r$, and an upper choice
gives $q_n=n+U(n-1)+1$ by Lemma~\ref{lem:upper}. The first thirty
rows are stored as a fixed seed.

\begin{algorithm}[!ht]
\caption{Handing on the next symbol of the queen word}
\label{alg:next-symbol}
\renewcommand{\algorithmicrequire}{\textbf{Input:}}
\renewcommand{\algorithmicensure}{\textbf{Output:}}
\begin{algorithmic}[1]
\Require A copy $C$: records \eqref{eq:generator-records} and a list of computed symbols.
\Ensure The next symbol that $C$ has not yet handed on, starting with $\sigma_{30}$.
\Function{NextSymbol}{$C$}
  \While{the list of $C$ is empty}
    \If{the next step of $C$ needs a symbol beyond the end of its $Q$}\Comment{a request}
      \If{$C$ has no input copy}
        \State Give $C$ an input copy, started before column $30$.
      \EndIf
      \State Append \Call{NextSymbol}{input copy of $C$} to the queue of $C$.
    \Else
      \State Carry out the step, and add its symbol $\sigma_n$ to the list.
    \EndIf
  \EndWhile
  \State Remove the first symbol from the list and return it.
\EndFunction
\end{algorithmic}
\end{algorithm}

Every symbol handed on is correct. Consider all steps of all copies
in the order in which they are carried out. The inputs of a step were
handed on earlier, so by induction they are actual symbols, and a
copy that receives actual symbols carries out the actual step, by the
induction of Section~\ref{sec:identification}. The simplifications
above change neither the queen chosen nor the records kept.

Every call also returns. Suppose that a copy $C$ before column $n$
calls its input copy $C'$, and that $C'$ has no symbol in its list.
The request is for $\sigma_k$ with $k=m+|Q|\le m+6<n$. So far $C'$
has handed on exactly $\sigma_{30},\dots,\sigma_{k-1}$, and with an
empty list it has computed exactly these, so it stands before
column $k<n$. Thus, by strong induction on the column, the call to
$C'$ returns: each step of $C'$ makes at most six requests, each to a
copy that either has a symbol ready or stands before a still smaller
column, and a copy before a column less than
$48$ needs no input at all.

\subsection{Compiling the calculation into a table}\label{sec:block-transducer}

In Algorithm~\ref{alg:next-symbol}, nearly all the work of a copy is
the calculation itself: the attack tests and record updates for every
queen. Along the actual process, however, the records take only
finitely many values. This section computes all of them in advance and
replaces the calculation by table lookups that place several queens
at once.

\medskip\noindent
\textbf{Paused records.}
Say that a copy is \emph{paused} when its next step needs one more
symbol in $Q$. Appending one input symbol lets the calculation run
until it pauses again, possibly placing several queens along the
way. Record all these outputs: for each queen, its symbol, its row
advance $\mu$, and its lower offset $r$ when it is lower. Thus one
input gives a new paused record and a list of outputs. The directed
graph of these operations is a finite \emph{transducer}: each edge
reads one symbol and produces a list of outputs.

For example, a copy first pauses before column $48$, after producing
$\sigma_{30},\dots,\sigma_{47}$, with $m=29$ and
\[
 w=1,\quad z=2,\quad R=D=A=\varnothing,\quad
 (u_{25},u_{26},u_{27},u_{28})=(0,0,1,0),\quad Q=\sigma_{29}=\mathtt{2}.
\]
The input $\sigma_{30}=\mathtt{1}$ is not yet enough, and the copy
pauses again with $Q=\mathtt{21}$ without placing a queen. The next
input, $\sigma_{31}=\mathtt{2}$, lets it place the lower queen
$(48,29)$. The input after that, $\sigma_{32}=\mathtt{2}$, places four
queens: the lower queen $(49,31)$ and the upper queens $(50,81)$,
$(51,83)$, and $(52,85)$. An input may thus yield no queen, one queen,
or several.

\medskip\noindent
\textbf{Which records occur.}
The table must cover every input that the actual process can supply.
To find these inputs, pair each paused record with the twelve symbols
ending at the right end of $Q$, which form a history-graph vertex.
Follow every outgoing edge of this vertex: append its symbol, compute
until the next pause, and shift the twelve symbols by one. Explore
all pairs reachable from the first pause in this way. Since the queen
word follows the history graph (Section~\ref{sec:identification}),
every actual input sequence is among those explored.

The history only decides which inputs to explore; the calculation
itself reads only the record and the supplied symbol. The exploration
reaches $2489$ pairs but only $300$ distinct paused records. We
therefore discard the histories, leaving a graph of $300$ paused
records whose $476$ edges keep their input symbols and complete
output lists.

\medskip\noindent
\textbf{Combining records and input steps.}
Some of these records have compatible responses and can share a table
position. We group them into $82$ classes, checking every defined
transition: whenever two records in a class accept the same input,
they must give identical output lists and successors in the same
class. An input undefined at a record imposes no condition there. The
program still obtains its input from the next copy; grouping records
does not determine that input.

We then combine four consecutive input transitions into one table
entry. These paths are formed on the $300$-record graph before
replacing their endpoints by classes. This order matters: paths
introduced only by merging have not been checked against the local
calculation. Whenever paths begin in the same class and read the same
four symbols, we check that their complete output lists agree and
that they end in the same class. Hence the class and the four input
symbols, packed into one byte, determine a single entry. Each entry
emits between three and twelve symbols. Algorithm~\ref{alg:table}
summarizes the construction.

\begin{algorithm}[!ht]
\caption{Building the table}
\label{alg:table}
\renewcommand{\algorithmicensure}{\textbf{Output:}}
\begin{algorithmic}[1]
\Ensure The table $T$, or failure.
\State $S_0\gets$ the records before column $30$, run until they pause; $W_0\gets\sigma_{18}\dots\sigma_{29}$.
\State $\mathcal P\gets\{(S_0,W_0)\}$, unexamined.
\While{some pair $(S,W)\in\mathcal P$ is unexamined}
  \State Mark $(S,W)$ as examined.
  \For{each edge $W\xrightarrow{s}W'$ of the history graph}
    \State Append $s$ to the queue of $S$ and run until the next pause at $S'$, with outputs $L$.
    \State Record the transition $S\to S'$ reading $s$ and emitting $L$.\Comment{fails if it differs}
    \State Add $(S',W')$ to $\mathcal P$ as unexamined, unless it is already in $\mathcal P$.
  \EndFor
\EndWhile
\State Group the paused records into compatible classes.
\For{each path of four transitions from $S$ to $S'$, reading $a_1a_2a_3a_4$}
  \State $T[\text{class of }S,\ a_1a_2a_3a_4]\gets(\text{joined outputs},\ \text{class of }S')$\Comment{fails if it differs}
\EndFor
\State \Return the table $T$
\end{algorithmic}
\end{algorithm}

\medskip\noindent
\textbf{Using the table.}
At run time a copy keeps only its table position, the class of its
current paused record. Given four input symbols packed into a byte,
it reads the entry at its position and that byte, emits the entry's
output symbols, and moves to the entry's successor position. One
lookup replaces four input steps of Algorithm~\ref{alg:next-symbol},
with all their attack tests and record updates.

The entries are stored in a fixed array of $64$-bit words. Each class
has a base offset into this array: adding the input byte gives the
location to read. The entry contains the successor's base offset, the
packed output symbols, their number, and a reference to their
coordinate data. The constructor checks every defined entry and
ensures that their storage locations do not conflict. During
generation the program retains this base offset as its table
position; neither graph nor any attack record is consulted. The
README of the generator gives the full construction counts and
storage layout.

The coordinate data are relative to the counters at the beginning of
the block. Suppose that the block starts at column $n_0$, with row
reference $m_0$ and upper count $U_0=U(n_0-1)$. Number its placements
starting at $0$, let $\mu_i$ be the row advance after placement $i$,
and let $r_j$ be the lower offset at placement $j$ when it is lower.
Summing the counter updates gives
\begin{equation}\label{eq:block-coordinates}
 q_{n_0+j}=
 \begin{cases}
  m_0+\displaystyle\sum_{i<j}\mu_i+r_j,
    &\text{if the queen is lower},\\[4pt]
  n_0+U_0+j+\displaystyle\sum_{i=0}^{j}u_{n_0+i},
    &\text{if the queen is upper}.
 \end{cases}
\end{equation}
For each row, the table therefore stores a small offset from one of
the two bases $m_0$ and $n_0+U_0$. It also stores the total advances
of $m$ and $U$ across the block. This gives exact integer coordinates
without repeating the attack tests. The outer copy can stop partway
through a block and resume without repeating either a lookup or a
counter update. The table, coordinate data, and thirty seed rows
occupy $34925$ bytes in total.

\subsection{Buffers, correctness, and cost}\label{sec:generator-cost}

With the table, the copies pass symbols four at a time, packed into
one byte, the input of a lookup. The remaining overhead lies mostly
in the calls between copies, so each copy also computes ahead: it
keeps a buffer of bytes ready for the copy on its left and refills it
in batches. The outermost input copy has a buffer budget of $1024$
bytes; each successive budget is halved, down to a minimum of one.
Each copy starts with the eighteen symbols
$\sigma_{30},\ldots,\sigma_{47}$: four complete bytes and two symbols
carried forward. Three padding bytes per buffer accommodate this
initial prefix and a block that goes past the budget.

Algorithm~\ref{alg:recursive-generation} gives the refill rule. It is
Algorithm~\ref{alg:next-symbol} with bytes in place of symbols, a
table lookup in place of four input steps, and a buffer refilled in
batches. Each copy
saves its table position, its place in the buffer, and up to three
symbols not yet forming a byte. The outer coordinate copy uses the
same table entries but decodes \eqref{eq:block-coordinates} instead
of buffering the output word.

\begin{algorithm}[!ht]
\caption{Supplying four earlier symbols to a table lookup}
\label{alg:recursive-generation}
\renewcommand{\algorithmicrequire}{\textbf{Input:}}
\renewcommand{\algorithmicensure}{\textbf{Output:}}
\begin{algorithmic}[1]
\Require A persistent copy $C$, initialized with the eighteen-symbol prefix.
\Ensure The next four symbols of its output, packed into one byte.
\Function{NextByte}{$C$}
  \If{the ready-byte buffer of $C$ is empty}
    \If{$C$ has no input copy}
      \State Create $C'$ at the same seed, with half the budget of $C$, at least one.
      \State Save $C'$ as the input copy of $C$.
    \EndIf
    \State Let $C'$ be the input copy of $C$.
    \Repeat
      \State $a\gets\Call{NextByte}{C'}$
      \State Look up the transition for the saved table position of $C$ and byte $a$.
      \State Save its successor table position and append its output symbols.
      \State Move complete groups of four symbols to the ready-byte buffer.
    \Until{the buffer contains at least the budget of $C$ bytes}
  \EndIf
  \State \Return the next ready byte, advancing the buffer position.
\EndFunction
\end{algorithmic}
\end{algorithm}

\medskip\noindent
\textbf{Correctness of the outputs.}
The finite table checks establish that correct input symbols give the
same outputs and updates as the calculation. They cover every
reachable pair of a history and a local record, every transition
after merging, and every four-input path. The argument of
Section~\ref{sec:generator-records} then applies unchanged: in the
order in which outputs are completed across all copies, each output
uses input already produced by the next copy, so by induction its
table entry gives the correct new symbols and updates, and buffering
preserves their order. A copy may compute ahead of the queen
currently being decoded by the outermost copy; the argument uses the
order of computation, rather than an ordering of the copies' current
queen indices.

\medskip\noindent
\textbf{Termination of requests.}
In Algorithm~\ref{alg:next-symbol}, a copy that must compute stands
before a smaller column than the copy that called it. With buffering,
a copy may have computed ahead of what it has handed on, so we count
complete bytes instead of columns. The key fact is that each copy has
computed more complete output bytes than it has consumed input bytes.
To include the effects of buffering, count all computed symbols,
including those still waiting to be supplied to another copy.

At a block boundary, let $t$ be the number of input bytes consumed,
let $H$ be the number of output symbols computed starting at
$\sigma_{30}$, and put $P=\lfloor H/4\rfloor$. Thus $H$ includes the
initial eighteen symbols, and $P$ counts complete output bytes,
whether already supplied or still buffered. The input endpoint and the
current output column satisfy
\[
 m+|Q|=30+4t,\qquad n=30+H=m+\kappa+z.
\]
At the initial pause, after computing through column $47$, the local
calculation has $n=48$, $m=29$, and $\kappa=U(28)=17$. The table
construction checks $z-|Q|\ge-4$ for all $300$ paused records, and
$\kappa$ never decreases along the actual process. Subtracting the
two identities therefore gives
\begin{equation}\label{eq:generator-lead}
 H-4t=\kappa+z-|Q|\ge13,
 \qquad P\ge t+3.
\end{equation}
In particular, the copy has computed at least three more complete
bytes than it has consumed. These counters describe the local
calculation represented by the table, even though the inner copies do
not store them explicitly.

First consider the part of the chain where every buffer budget is one
byte. A refill requests more input only while it has produced no
complete new byte; as soon as it has one, it returns. If it calls for
a refill of its input copy, that copy's buffer is exhausted. The
input copy has therefore computed exactly the $t$ complete bytes
consumed by the requesting copy. By \eqref{eq:generator-lead}, its
count is at least three less than the requesting copy's $P$. Thus
every further recursive call within this refill has a strictly
smaller nonnegative count.

Strong induction on $P$ now proves termination. A fresh copy already
has four ready bytes. Otherwise each deeper request returns by the
induction hypothesis, and each successful table lookup adds at least
three symbols. At most two lookups therefore complete a new byte and
finish the refill. Only finitely many copies precede this part of the
chain with budgets greater than one, and each of their refills needs
finitely many requests to its input copy. Hence every request
returns.

\medskip\noindent
\textbf{Time and space bounds.}
The strict decrease above proves termination. To obtain the stronger
bounds on work and storage, we make precise the estimate of
Section~\ref{sec:generator-idea}: successive copies compute prefixes
whose lengths decrease geometrically, apart from a bounded buffering
allowance.

\begin{proposition}[Sequential generation]\label{prop:stream-generation}
With a fixed initial buffer budget, the table algorithm generates
$q_0,\ldots,q_N$ exactly in $O(N)$ operations on $O(\log N)$-bit
machine words. Its working memory is $O(\log N)$ words, including
the recursive stack, in addition to a fixed table and excluding
retained output.
\end{proposition}
\begin{proof}
Correctness and termination were established above. For the cost,
Proposition~\ref{prop:stability} gives $\kappa=m/\phi+O(1)$.
Substituting this into the identities for the input endpoint and
output column, with bounded $z$ and $|Q|$, gives $t=P/\phi+O(1)$.
Thus a copy computing $P$ output bytes consumes only $P/\phi+O(1)$
input bytes.

Let $P_j$ count complete bytes computed by the $j$th copy, and let
$b_{j+1}$ be the next copy's buffer budget. Between refills, the next
copy can have computed more bytes than were requested only by the
size of its ready buffer and a bounded amount of padding.
Consequently
\begin{equation}\label{eq:generator-contraction}
 P_{j+1}\le\frac{P_j}{\phi}+O(b_{j+1}+1).
\end{equation}
With a fixed initial budget, this gives geometric decrease above a
fixed threshold. Only a fixed number of levels have budgets greater
than one. Below the threshold, the strict decrease in
\eqref{eq:generator-lead} bounds the remaining active chain of
budget-one copies. A new copy is allocated only along an active
chain, so at most $O(\log N)$ copies are ever retained.

Summing \eqref{eq:generator-contraction} along the chain gives $O(N)$
total byte production and hence $O(N)$ table lookups. Each lookup,
byte transfer, and coordinate decoding has bounded cost. For $L$
inner copies with initial budget $B$, the halving budgets sum to at
most $2B+L$, and the padding uses $3L$ further bytes. Each copy also
retains a bounded record and a bounded stack frame. Since $B$ is
fixed and $L=O(\log N)$, the claimed space bound follows.
\end{proof}

Knuth's \texttt{infty-queens} program \cite{KnuthQueens} retains
occupancy arrays growing linearly with the requested number of
queens. Here the fixed table and short chain replace those arrays.
The algorithm is still sequential: computing $q_N$ entails computing
the preceding queens. A repeated local record does not determine its
future inputs, so the finite table alone does not provide random
access to a distant queen.

\subsection{C implementation and measured performance}\label{sec:generation-benchmarks}

On a fresh $64$~GiB Ubuntu EC2 instance, our C generator produced and
hashed ten billion queen rows in a median $25.413$~s, compared with
$49.040$~s for bit-packed Knuth: a speedup of $1.93$. Its median peak
memory use was $1.76$~MiB, compared with $6243.51$~MiB.

\begin{table}[!ht]
\centering
\begin{tabular}{@{}r r r r r@{}}
\toprule
& \multicolumn{2}{c}{Elapsed time (s)}
& \multicolumn{2}{c}{Peak memory (MiB)}\\
\cmidrule(lr){2-3}\cmidrule(l){4-5}
$M$ & Our C & Packed Knuth & Our C & Packed Knuth\\
\midrule
$10^{6}$ & 0.003 & 0.005 & 1.76 & 2.20\\
$10^{7}$ & 0.026 & 0.049 & 1.67 & 7.82\\
$10^{8}$ & 0.253 & 0.486 & 1.76 & 64.07\\
$10^{9}$ & 2.524 & 4.862 & 1.67 & 625.88\\
$10^{10}$ & 25.413 & 49.040 & 1.76 & 6243.51\\
\bottomrule
\end{tabular}
\caption{Measured generation and hashing on the Ubuntu EC2
instance. Both columns use C, with the same checksum and
measurement boundary.}
\label{tab:c-generation}
\end{table}

The implementation in the folder \path{fast_generator} of the
companion repository uses portable
C11 with fixed-width unsigned integers. Its default table reads four
symbols at a time, with initial buffer budget $1024$, and it loads no
graph or data files at run time. The interface supports single rows,
filled arrays, and incremental hashing, with arbitrary pauses between
calls.

We compared it with a bit-packed adaptation of Knuth's
\texttt{infty-queens} \cite{KnuthQueens}. The adaptation packs the
occupancy flags into $64$-bit words and widens indices to $64$ bits,
retaining Knuth's placement order and control flow; its statistics
and per-queen printing are removed. Both programs decode every row,
fold it into the same $64$-bit checksum, and print one final summary
without retaining the sequence. Before timing, all table artifacts
were regenerated, all $3483$ four-input paths, comprising $24811$
local steps, were replayed against a separate implementation of the
calculation, and every coordinate through the first million queens was
compared with an independent occupancy-based calculation. These finite
tests support the implementation; its correctness is
Proposition~\ref{prop:stream-generation}.

Table~\ref{tab:c-generation} reports medians of three runs after one
excluded warmup, on an Amazon EC2 \texttt{r7i.2xlarge} instance with
an Intel Xeon Platinum 8488C processor, Ubuntu~24.04, and GCC~13.3.
Both programs used the same compiler options, ran single-threaded,
and were pinned to one CPU, alternating in order. Elapsed time runs
from process creation to exit. Peak memory is the largest amount of
physical memory that the process occupied at any time (its peak
resident set size), as reported by the kernel. All runs agreed on the final coordinate and checksum. The
count $M$ includes the origin, so the last column is $M-1$.

From $10^9$ to $10^{10}$ queens, both elapsed times grew by a factor
of about ten, consistent with linear running time. The generator's
fixed data occupy $34925$ bytes, and at $10^{10}$ queens its $42$
levels request a further $3560$ bytes of heap. This storage grows
logarithmically, although at these sizes it fits within memory pages
that the process has already allocated. Knuth's packed arrays instead
need about $(2\phi+2)M/8\approx0.655M$ bytes: about $61$~GiB at
$10^{11}$ queens, close to the capacity of this machine. These results
concern bulk generation with a checksum on this machine; they do not
establish the same ratio for row-at-a-time calls, printed output, or
other hardware. The report \path{fast_generator/REPORT.md} gives the full measurement
procedure, raw observations, and linear projections to larger counts,
and Appendix~\ref{app:stream-generation} gives the commands.
